\section{Introducción: ¿cuál es el hilo conductor de este informe trimestral?}

Esta sección establece primero el marco de lectura de todo el episodio. El «Informe trimestral global de grandes modelos» de Zhang Xiaojun (张小珺) no es un recuento de noticias una por una, sino que usa un conjunto de juicios para condensar los cambios del trimestre en una hoja de ruta tecnológica. El cambio clave del primer trimestre de 2025 es que Li Guangmi (李广密) vuelve a subrayar el papel fundamental del Pre-training y, al mismo tiempo, discute Coding, Agent, Online Learning, la divergencia estratégica de las empresas de modelos y el panorama entre China y Estados Unidos dentro de un mismo hilo conductor hacia la AGI.

\begin{importantbox}{Tesis central de este episodio}
El juicio central de Guangmi puede condensarse en una sola frase: el aumento de la inteligencia sigue siendo el único hilo conductor; el Pre-training se encarga de abrir el límite superior intrínseco del modelo, Coding ofrece el entorno de acción digital más general, Agent es una nueva especie y Online Learning podría ser la siguiente ruta de nivel paradigmático; tanto las empresas de producto como las empresas de modelos deben reordenar su posición en torno a estas capacidades.
\end{importantbox}

\begin{knowledgebox}{Nota sobre la estrategia visual}
Este video es una toma fija de pódcast, sin proyección de pantalla, pizarrón ni demostración de producto. En el texto, la portada se usa solo para identificar la fuente; todas las imágenes del texto son diagramas conceptuales propios, que sirven para explicar rutas tecnológicas, estrategias organizacionales, barreras de producto y el panorama de la industria. Esto se ajusta mejor al criterio de este repositorio para el tratamiento de pódcasts que repetir tomas del conductor o del invitado.
\end{knowledgebox}

\begin{figure}[H]
\centering
\includegraphics[width=0.92\textwidth]{q1-model-map.png}
\caption{Mapa del informe trimestral de grandes modelos del primer trimestre de 2025: Pre-training, Coding, Agent, Online Learning y el panorama global conforman juntos el hilo conductor. Diagrama conceptual propio, elaborado a partir de 00:00:08--00:04:22 y del contenido de todo el video.}
\end{figure}

\begin{knowledgebox}{Lectura de la figura: no leer el informe trimestral como una lista de noticias}
En la figura, Q1 2025 está en el centro, rodeado por Pre-training, Coding, Agent, Online, producto y panorama. Al leer esta figura, observe primero las relaciones de dependencia: la capacidad del modelo determina el límite superior utilizable, Coding y Agent ofrecen el entorno de acción, los productos recogen los beneficios de la inteligencia y el panorama de las empresas refleja cómo apuestan las organizaciones por estas rutas.
\end{knowledgebox}

\subsection{Ruta de lectura}

Esta sección presenta la ruta de lectura. Todo el video puede desglosarse en seis preguntas: primera, ¿por qué, cuando todos se orientan hacia las aplicaciones y el post-entrenamiento, Guangmi vuelve a subrayar el Pre-training? Segunda, ¿por qué Coding no es solo programación, sino el entorno más importante en la etapa temprana de la AGI? Tercera, ¿qué problema organizacional revela la divergencia estratégica entre OpenAI y Anthropic? Cuarta, ¿por qué se dice que Agent es una nueva especie? Quinta, ¿puede Online Learning cambiar el paradigma de entrenamiento? Sexta, ¿cómo se restringen mutuamente los modelos, los productos, el capital y el panorama entre China y Estados Unidos?

\noindent\begin{tabularx}{\textwidth}{p{0.19\textwidth}p{0.34\textwidth}X}
\toprule
Pregunta de lectura & Material de la entrevista & Juicio que hay que formar \\
\midrule
¿De dónde viene el límite superior del modelo? & Pre-training, Post-training, RL, datos sintéticos & El post-entrenamiento puede dar forma, pero el límite superior de la base sigue proviniendo del preentrenamiento y de los datos, el cómputo y la arquitectura. \\
¿Cómo se lleva Agent a la práctica? & Coding, computer use, tool use, long context & El entorno digital ofrece retroalimentación verificable; es donde Agent llega primero a la práctica. \\
¿Por qué divergen las empresas? & OpenAI, Anthropic, DeepSeek, Manus, Cursor & La estrategia no es un eslogan, sino la expresión de las capacidades organizacionales, los puntos de entrada de producto y la acumulación tecnológica. \\
¿Dónde está la barrera del producto? & Cloud, OS, contenedores, costo de token, «ladrones del fuego» & En la era del modelo desnudo se debilita; se fortalecen los entornos y flujos de trabajo que recogen la inteligencia. \\
¿Cuál es el siguiente paradigma? & Online Learning, long-term memory, retroalimentación del entorno & Si el modelo puede aprender mientras actúa, la ruta de aumento de la inteligencia volverá a cambiar. \\
\bottomrule
\end{tabularx}

\begin{warningbox}{Límite importante: este es un informe trimestral de opinión, no un anuario de hechos}
Esta nota organiza las opiniones del invitado en un marco de análisis, pero no presenta como hechos juicios del tipo «la AGI se logrará en dos años», «los pesos de la cartera de cierta empresa» o «cierta empresa fracasará». Otro límite terminológico: el Perplexity que aparece en este episodio se refiere al producto o empresa de búsqueda con IA, no a la perplexity de la evaluación de modelos de lenguaje (PPL, el indicador de perplejidad que se obtiene al aplicar la exponencial a la entropía cruzada y que intuitivamente puede entenderse como el número efectivo de candidatos por token).
\end{warningbox}

\subsection{Resumen de la sección}

El valor técnico de este informe trimestral radica en que reorganiza varias líneas de interés del primer trimestre de 2025 en un solo hilo conductor hacia la AGI: el límite superior del modelo, el entorno de acción digital, la conversión de Agent en producto, el paradigma de aprendizaje en línea, la estrategia de las empresas y las restricciones geopolíticas. Las secciones siguientes lo desarrollan nivel por nivel.
